% GENERATED FILE. Edit canonical lessons, then run npm run generate:books.
% canonical-source-hash: fnv1a64:58a0e45c0a463ab3
% canonical-lessons: PT-C517-flauta, PT-C517-trompete, PT-C517-saxofone, PT-C517-acordeao, PT-C517-afinar

\chapter{Flautas e trompetes}
\label{ch:pt-flautas-e-trompetes}
\hlchaptermodality{517}

\begin{chapteropening}
\textbf{By the end of this chapter:} \emph{I can talk about wind instruments and tuning.}

Complete the last lesson of chapter 517: I can talk about wind instruments and tuning.
\end{chapteropening}

\section[a flauta]{a flauta — flute; recorder}

\label{lesson:PT-C517-flauta}



\begin{quote}
Before the new one: say the Portuguese for cello, then the Portuguese for harp.
\end{quote}

\subsection*{You'll want to know: a flauta}
\textbf{a flauta} — \textquotedblleft{}flute; recorder\textquotedblright{}.

\begin{quote}
Na escola, todos aprendem a tocar flauta.
\end{quote}

\begin{quote}
\emph{At school, everyone learns to play the recorder.}
\end{quote}

\emph{Flauta de bisel} — recorder; \emph{flauta transversal} — flute.

\begin{cousinweb}
\emph{flauta} is usually traced to Old Occitan \emph{flaut}, a word usually taken to imitate the sound of blowing; the origin is not settled.
\end{cousinweb}

\begin{grammarlens}[title={What you've built}]
One more word for telling, explaining and describing at B1.
\end{grammarlens}

\subsection*{Your turn}
\begin{itemize}
\raggedright
  \item \emph{Say it:} \emph{a flauta}
  \item \emph{Say it:} \emph{a flauta}, once more
  \item \emph{Say it:} a sentence of your own with \emph{a flauta}
  \item \emph{Recall:} say the Portuguese for cello, then the Portuguese for harp, then say \emph{a flauta} again
\end{itemize}

\begin{tcolorbox}[breakable,colback=teal!4,colframe=teal!35!black,title={Before you move on}]
What does a flauta mean? (\textquotedblleft{}Flute; recorder\textquotedblright{}.) Say it once more.
\end{tcolorbox}
\section[o trompete]{o trompete — trumpet}

\label{lesson:PT-C517-trompete}



\begin{quote}
Before the new one: say the Portuguese for harp, then the Portuguese for flute.
\end{quote}

\subsection*{You'll want to know: o trompete}
\textbf{o trompete} — \textquotedblleft{}trumpet\textquotedblright{}.

\begin{quote}
O vizinho toca trompete à noite.
\end{quote}

\begin{quote}
\emph{The neighbour plays the trumpet at night.}
\end{quote}

\emph{A trombeta} is an older word for a trumpet.

\begin{cousinweb}
\emph{trompete} is usually traced to French \emph{trompette}, a diminutive of \emph{trompe} \textquotedblleft{}a horn\textquotedblright{}, which goes back to a Germanic word.
\end{cousinweb}

\begin{grammarlens}[title={What you've built}]
One more word for telling, explaining and describing at B1.
\end{grammarlens}

\subsection*{Your turn}
\begin{itemize}
\raggedright
  \item \emph{Say it:} \emph{o trompete}
  \item \emph{Say it:} \emph{o trompete}, once more
  \item \emph{Say it:} a sentence of your own with \emph{o trompete}
  \item \emph{Recall:} say the Portuguese for harp, then the Portuguese for flute, then say \emph{o trompete} again
\end{itemize}

\begin{tcolorbox}[breakable,colback=teal!4,colframe=teal!35!black,title={Before you move on}]
What does o trompete mean? (\textquotedblleft{}Trumpet\textquotedblright{}.) Say it once more.
\end{tcolorbox}
\section[o saxofone]{o saxofone — saxophone}

\label{lesson:PT-C517-saxofone}



\begin{quote}
Before the new one: say the Portuguese for flute, then the Portuguese for trumpet.
\end{quote}

\subsection*{You'll want to know: o saxofone}
\textbf{o saxofone} — \textquotedblleft{}saxophone\textquotedblright{}.

\begin{quote}
Ouvimos um solo de saxofone no bar.
\end{quote}

\begin{quote}
\emph{We heard a saxophone solo in the bar.}
\end{quote}

Often shortened to \emph{o sax}.

\begin{cousinweb}
\emph{saxofone} comes from French \emph{saxophone} (EtymDB), built on the surname of its Belgian inventor and Greek \emph{phōn\'{\={e}}} \textquotedblleft{}sound\textquotedblright{}.
\end{cousinweb}

\begin{grammarlens}[title={What you've built}]
One more word for telling, explaining and describing at B1.
\end{grammarlens}

\subsection*{Your turn}
\begin{itemize}
\raggedright
  \item \emph{Say it:} \emph{o saxofone}
  \item \emph{Say it:} \emph{o saxofone}, once more
  \item \emph{Say it:} a sentence of your own with \emph{o saxofone}
  \item \emph{Recall:} say the Portuguese for flute, then the Portuguese for trumpet, then say \emph{o saxofone} again
\end{itemize}

\begin{tcolorbox}[breakable,colback=teal!4,colframe=teal!35!black,title={Before you move on}]
What does o saxofone mean? (\textquotedblleft{}Saxophone\textquotedblright{}.) Say it once more.
\end{tcolorbox}
\section[o acordeão]{o acordeão — accordion}

\label{lesson:PT-C517-acordeao}



\begin{quote}
Before the new one: say the Portuguese for trumpet, then the Portuguese for saxophone.
\end{quote}

\subsection*{You'll want to know: o acordeão}
\textbf{o acordeão} — \textquotedblleft{}accordion\textquotedblright{}.

\begin{quote}
Nas festas da aldeia, o avô tocava acordeão.
\end{quote}

\begin{quote}
\emph{At the village fêtes, Grandpa played the accordion.}
\end{quote}

In Brazil: \emph{o acordeom} or \emph{a sanfona}.

\begin{cousinweb}
\emph{acordeão} was borrowed from French \emph{accordéon} (EtymDB), from German \emph{Akkordion}, built on \emph{Akkord} \textquotedblleft{}a chord\textquotedblright{}.
\end{cousinweb}

\begin{grammarlens}[title={What you've built}]
One more word for telling, explaining and describing at B1.
\end{grammarlens}

\subsection*{Your turn}
\begin{itemize}
\raggedright
  \item \emph{Say it:} \emph{o acordeão}
  \item \emph{Say it:} \emph{o acordeão}, once more
  \item \emph{Say it:} a sentence of your own with \emph{o acordeão}
  \item \emph{Recall:} say the Portuguese for trumpet, then the Portuguese for saxophone, then say \emph{o acordeão} again
\end{itemize}

\begin{tcolorbox}[breakable,colback=teal!4,colframe=teal!35!black,title={Before you move on}]
What does o acordeão mean? (\textquotedblleft{}Accordion\textquotedblright{}.) Say it once more.
\end{tcolorbox}
\section[afinar]{afinar — to tune}

\label{lesson:PT-C517-afinar}



\begin{quote}
Before the new one: say the Portuguese for saxophone, then the Portuguese for accordion.
\end{quote}

\subsection*{You'll want to know: afinar}
\textbf{afinar} — \textquotedblleft{}to tune\textquotedblright{}.

\begin{quote}
Antes do concerto, os músicos afinam os instrumentos.
\end{quote}

\begin{quote}
\emph{Before the concert, the musicians tune their instruments.}
\end{quote}

\emph{Desafinar} — to sing or play out of tune.

\begin{cousinweb}
\emph{afinar} is built from \emph{a-} and \emph{fino} \textquotedblleft{}fine, thin\textquotedblright{}: to make a sound fine and exact. \emph{Fino} is usually traced to Latin \emph{fīnis} \textquotedblleft{}end\textquotedblright{} (EtymDB), but the path of meaning is not settled.
\end{cousinweb}

\begin{grammarlens}[title={What you've built}]
One more word for telling, explaining and describing at B1.
\end{grammarlens}

\subsection*{Your turn}
\begin{itemize}
\raggedright
  \item \emph{Say it:} \emph{afinar}
  \item \emph{Say it:} \emph{afinar}, once more
  \item \emph{Say it:} a sentence of your own with \emph{afinar}
  \item \emph{Recall:} say the Portuguese for saxophone, then the Portuguese for accordion, then say \emph{afinar} again
\end{itemize}

\begin{tcolorbox}[breakable,colback=teal!4,colframe=teal!35!black,title={Before you move on}]
What does afinar mean? (\textquotedblleft{}To tune\textquotedblright{}.) Say it once more.
\end{tcolorbox}
